% ============================================================
% 分节契约
% purpose:      结果意味着什么、边界在哪、后续方向
% allowed claims: 从主结果直接推出的结论;局限须与范围边界一致
% forbidden claims: 超出证据的推广;防御性自我攻击(有礼有节地讲局限)
% inputs:       argument_map.md · 01-topic.md 范围边界
% 验证方式:     每个结论指向对应结果;anti-defensive-writing 过立场
% ============================================================
\section{Discussion}\label{sec:discussion}

<第 1 段:主结果回答了引言提出的问题到什么程度。>

<第 2 段:局限——与 topic.md 的范围边界一致,陈述事实不示弱。>

<第 3 段:后续方向。>

% Statistical analysis 小节(nature-statistics 审查后落笔):
% \subsection{Statistical analysis}
% <重复数、检验方法、多重比较校正、置信区间。>
